% ============================================================
%  MTH636 -- Game Theory | Lecture 6
%  Two-Player Zero-Sum Games & Min-Max Theorem
% ============================================================
\documentclass[12pt, a4paper]{article}

% ============================================================
%  MTH636 -- Game Theory
%  Notes Preamble — Color-Coded, TikZ-Enhanced
% ============================================================

% --- Core packages ---
\usepackage[T1]{fontenc}
\usepackage[utf8]{inputenc}
\usepackage{lmodern}
\usepackage[margin=2.5cm, top=2.8cm, bottom=2.8cm]{geometry}
\usepackage{microtype}

% --- Math ---
\usepackage{amsmath, amssymb, amsthm, mathtools}
\usepackage{bm}        % bold math
\usepackage{cancel}    % strike-through in math

% --- Color & Boxes ---
\usepackage[dvipsnames, svgnames, x11names, table]{xcolor}
\usepackage[most]{tcolorbox}
\tcbuselibrary{breakable, skins, theorems}

% --- TikZ & Diagrams ---
\usepackage{tikz}
\usepackage{tikz-qtree}
\usetikzlibrary{arrows.meta, positioning, shapes.geometric, fit, backgrounds,
                decorations.pathreplacing, calc, matrix, trees}
\usepackage{pgfplots}
\pgfplotsset{compat=1.18}

% --- Tables ---
\usepackage{booktabs}
\usepackage{multirow}
\usepackage{array}
\usepackage{tabularx}
\usepackage{makecell}

% --- Typography & Layout ---
\usepackage{enumitem}
\usepackage{parskip}
\usepackage{fancyhdr}
\usepackage{titlesec}
\usepackage{titletoc}
\usepackage{hyperref}
\usepackage{graphicx}
\usepackage{float}
\usepackage{caption}

% --- Color Palette ---
\definecolor{defblue}{HTML}{1A5276}      % Definition header
\definecolor{defbluebg}{HTML}{D6EAF8}    % Definition background
\definecolor{thmgreen}{HTML}{1E8449}     % Theorem header
\definecolor{thmbg}{HTML}{D5F5E3}        % Theorem background
\definecolor{exorange}{HTML}{B7770D}     % Example header
\definecolor{exbg}{HTML}{FEF9E7}         % Example background
\definecolor{warnred}{HTML}{922B21}      % Warning/Key Insight header
\definecolor{warnbg}{HTML}{FADBD8}       % Warning background
\definecolor{noticegray}{HTML}{2C3E50}   % Note header
\definecolor{noticebg}{HTML}{EBF5FB}     % Note background
\definecolor{taketawayyellow}{HTML}{7D6608} % Key Takeaway header
\definecolor{takeawaybg}{HTML}{FDFEFE}   % Takeaway background
\definecolor{accentpurple}{HTML}{6C3483} % Accent / proof
\definecolor{lightgray}{HTML}{F2F3F4}    % shading

% ============================================================
%  TCOLORBOX ENVIRONMENTS
% ============================================================

% -- Definition (Blue) --
\newtcolorbox{defbox}[2][]{
  enhanced, breakable,
  colback=defbluebg, colframe=defblue, coltitle=white,
  fonttitle=\bfseries\sffamily,
  title={Definition: #2},
  attach boxed title to top left={yshift=-2mm, xshift=4mm},
  boxed title style={colback=defblue, rounded corners},
  arc=4pt, boxrule=1.2pt,
  left=6pt, right=6pt, top=8pt, bottom=6pt,
  #1
}

% -- Theorem / Proposition (Green) --
\newtcolorbox{thmbox}[2][]{
  enhanced, breakable,
  colback=thmbg, colframe=thmgreen, coltitle=white,
  fonttitle=\bfseries\sffamily,
  title={Theorem: #2},
  attach boxed title to top left={yshift=-2mm, xshift=4mm},
  boxed title style={colback=thmgreen, rounded corners},
  arc=4pt, boxrule=1.2pt,
  left=6pt, right=6pt, top=8pt, bottom=6pt,
  #1
}

% -- Proposition (Green variant) --
\newtcolorbox{propbox}[2][]{
  enhanced, breakable,
  colback=thmbg, colframe=thmgreen!70!black, coltitle=white,
  fonttitle=\bfseries\sffamily,
  title={Proposition: #2},
  attach boxed title to top left={yshift=-2mm, xshift=4mm},
  boxed title style={colback=thmgreen!70!black, rounded corners},
  arc=4pt, boxrule=1.2pt,
  left=6pt, right=6pt, top=8pt, bottom=6pt,
  #1
}

% -- Example (Orange/Yellow) --
\newtcolorbox{exbox}[2][]{
  enhanced, breakable,
  colback=exbg, colframe=exorange, coltitle=white,
  fonttitle=\bfseries\sffamily,
  title={Example: #2},
  attach boxed title to top left={yshift=-2mm, xshift=4mm},
  boxed title style={colback=exorange, rounded corners},
  arc=4pt, boxrule=1.2pt,
  left=6pt, right=6pt, top=8pt, bottom=6pt,
  #1
}

% -- Key Insight / Warning (Red) --
\newtcolorbox{warnbox}[2][]{
  enhanced, breakable,
  colback=warnbg, colframe=warnred, coltitle=white,
  fonttitle=\bfseries\sffamily,
  title={Key Insight: #2},
  attach boxed title to top left={yshift=-2mm, xshift=4mm},
  boxed title style={colback=warnred, rounded corners},
  arc=4pt, boxrule=1.2pt,
  left=6pt, right=6pt, top=8pt, bottom=6pt,
  #1
}

% -- Note (Gray-Blue) --
\newtcolorbox{notebox}[1][]{
  enhanced, breakable,
  colback=noticebg, colframe=noticegray!60, coltitle=white,
  fonttitle=\bfseries\sffamily,
  title={Note},
  attach boxed title to top left={yshift=-2mm, xshift=4mm},
  boxed title style={colback=noticegray!60, rounded corners},
  arc=4pt, boxrule=0.8pt,
  left=6pt, right=6pt, top=8pt, bottom=6pt,
  #1
}

% -- Key Takeaway --
\newtcolorbox{takebox}[1][]{
  enhanced, breakable,
  colback=takeawaybg, colframe=taketawayyellow, coltitle=white,
  fonttitle=\bfseries\sffamily,
  title={Key Takeaway},
  attach boxed title to top left={yshift=-2mm, xshift=4mm},
  boxed title style={colback=taketawayyellow, rounded corners},
  arc=4pt, boxrule=1.2pt,
  left=6pt, right=6pt, top=8pt, bottom=6pt,
  #1
}

% -- Intuition --
\newtcolorbox{intubox}[1][]{
  enhanced, breakable,
  colback=accentpurple!8, colframe=accentpurple!60, coltitle=white,
  fonttitle=\bfseries\sffamily,
  title={Intuition},
  attach boxed title to top left={yshift=-2mm, xshift=4mm},
  boxed title style={colback=accentpurple!80, rounded corners},
  arc=4pt, boxrule=0.8pt,
  left=6pt, right=6pt, top=8pt, bottom=6pt,
  #1
}

% ============================================================
%  PAYOFF MATRIX HELPER
%  Usage: inside tabular, use \payoff{a}{b} for (a,b)
% ============================================================
\newcommand{\payoff}[2]{$(#1,\;#2)$}
\newcommand{\nash}[2]{\underline{$#1$},\underline{$#2$}}  % underline NE payoffs
\newcommand{\Nset}{\mathbb{N}}
\newcommand{\Rset}{\mathbb{R}}
\newcommand{\Sset}{\mathcal{S}}
\newcommand{\Ui}{u_i}

% ============================================================
%  SECTION FORMATTING
% ============================================================
\titleformat{\section}
  {\large\bfseries\sffamily\color{defblue}}
  {\thesection}{1em}{}[\titlerule]

\titleformat{\subsection}
  {\normalsize\bfseries\sffamily\color{thmgreen!80!black}}
  {\thesubsection}{1em}{}

\titleformat{\subsubsection}
  {\normalsize\itshape\sffamily\color{exorange!80!black}}
  {\thesubsubsection}{1em}{}

% ============================================================
%  HEADER / FOOTER
% ============================================================
\pagestyle{fancy}
\fancyhf{}
\fancyhead[L]{\sffamily\small\color{defblue} MTH636 -- Game Theory}
\fancyhead[R]{\sffamily\small\color{defblue} \leftmark}
\fancyfoot[C]{\sffamily\small\thepage}
\renewcommand{\headrulewidth}{0.4pt}

% ============================================================
%  HYPERREF SETUP
% ============================================================
\hypersetup{
  colorlinks=true,
  linkcolor=defblue,
  urlcolor=thmgreen,
  citecolor=accentpurple,
  bookmarksnumbered=true,
  pdftitle={MTH636 Game Theory Notes},
  pdfauthor={MTH636 Student Notes}
}

% ============================================================
%  MISC HELPERS
% ============================================================
\newcommand{\BR}{\text{BR}}        % Best Response
\newcommand{\NE}{\text{NE}}        % Nash Equilibrium
\newcommand{\SPNE}{\text{SPNE}}    % Subgame Perfect NE
\newcommand{\Dom}{\text{dom}}      % Dominance
\newcommand{\mixs}[1]{\sigma_{#1}} % Mixed strategy
\newcommand{\pures}[1]{s_{#1}}     % Pure strategy


% ============================================================
\begin{document}

% Title
\begin{center}
  {\LARGE\bfseries\sffamily\color{defblue} MTH636 -- Game Theory}\\[6pt]
  {\Large\sffamily Lecture 6: Two-Player Zero-Sum Games}\\[4pt]
\end{center}
\vspace{0.5cm}
\hrule height 1.5pt
\vspace{0.3cm}

% ============================================================
\section{Zero-Sum Games: Max-Min and Min-Max}

Recall that in a two-player zero-sum game, $u_1(s_1, s_2) + u_2(s_1, s_2) = 0$. Let $u = u_1$. Thus, $u_2 = -u$.
Player 1 wants to maximize $u$, while Player 2 wants to maximize $u_2 = -u$, which is equivalent to Player 2 wanting to \emph{minimize} $u$.

\begin{exbox}{Zero-Sum Matrix Notation}
  Consider a game where Player 1 chooses U or NU, and Player 2 chooses R or NR.
  Let the payoffs for Player 1 be:
  $u(U, R) = 3$, $u(NU, R) = -2$, $u(U, NR) = 1$, $u(NU, NR) = 3$.
  
  We represent this simply with Player 1's payoffs:
  \begin{center}
  \renewcommand{\arraystretch}{1.3}
  \begin{tabular}{c|cc}
    $P_1 \setminus P_2$ & R & NR \\
    \hline
    U  & \payoff{3}{-3} & \payoff{1}{-1} \\
    NU & \payoff{-2}{2} & \payoff{3}{-3} \\
  \end{tabular}
  \end{center}
\end{exbox}

\subsection{Player 1: Max-Min Value ($\underline{v}$)}
Player 1 guarantees themselves a payoff of:
\[
  \underline{v}_1 = \max_{s_1 \in S_1} \min_{s_2 \in S_2} u_1(s_1, s_2) = \max_{s_1} \min_{s_2} u(s_1, s_2)
\]
We denote this as $\underline{v}$. This is the \textbf{maximum amount Player 1 can guarantee to get}.

\subsection{Player 2: Min-Max Value ($\overline{v}$)}
Player 2 looks at their own max-min value:
\[
  \underline{v}_2 = \max_{s_2 \in S_2} \min_{s_1 \in S_1} u_2(s_1, s_2) = \max_{s_2} \min_{s_1} (-u(s_1, s_2))
\]
Because $\max(-f) = -\min(f)$, we can rewrite this as:
\[
  \underline{v}_2 = -\min_{s_2 \in S_2} \max_{s_1 \in S_1} u(s_1, s_2)
\]
Let $\overline{v} = \min_{s_2} \max_{s_1} u(s_1, s_2)$. This is called the \textbf{min-max value}.
For Player 2, $\overline{v}$ represents the \textbf{minimum amount Player 2 can guarantee to pay} to Player 1.
In this context, for Player 2, we call the max-min strategy the \textbf{min-max strategy}.

% ============================================================
\section{The Relationship Between $\underline{v}$ and $\overline{v}$}

\begin{thmbox}{Inequality of Max-Min and Min-Max}
  In general, for any game:
  \[
    \underline{v} \leq \overline{v}
  \]
  Meaning: (The max amount P1 can guarantee to get) $\leq$ (The min amount P2 can guarantee to pay).
\end{thmbox}

\begin{exbox}{Example 1: Value Exists}
  \begin{center}
  \renewcommand{\arraystretch}{1.3}
  \begin{tabular}{c|ccc|c}
    & L & C & R & \textbf{Row Min} \\
    \hline
    T & 3 & -5 & -2 & -5 \\
    M & 1 & 4 & 1 & \textbf{\color{defblue} 1} $\implies \underline{v} = 1$ \\
    B & 6 & -3 & -5 & -5 \\
    \hline
    \textbf{Col Max} & 6 & 4 & \textbf{\color{defblue} 1} & $\implies \overline{v} = 1$
  \end{tabular}
  \end{center}
  Here, $\underline{v} = 1$ and $\overline{v} = 1$. The max-min strategy for P1 is M, and the min-max strategy for P2 is R.
\end{exbox}

\begin{exbox}{Example 2: No Value in Pure Strategies}
  \begin{center}
  \renewcommand{\arraystretch}{1.3}
  \begin{tabular}{c|cc|c}
    & L & R & \textbf{Row Min} \\
    \hline
    T & -2 & 5 & -2 \\
    B & 3 & 0 & \textbf{\color{defblue} 0} $\implies \underline{v} = 0$ \\
    \hline
    \textbf{Col Max} & \textbf{\color{defblue} 3} & 5 & $\implies \overline{v} = 3$
  \end{tabular}
  \end{center}
  Here, $\underline{v} = 0$ and $\overline{v} = 3$. Since $\underline{v} < \overline{v}$, this game has \textbf{no value in pure strategies}.
\end{exbox}

% ============================================================
\section{Value of a Game and Optimal Strategies}

\begin{defbox}{Value of a Game}
  For a two-player zero-sum game, if $\underline{v} = \overline{v} = v$, then we say the game has a \textbf{value} $v$.
  
  And, the max-min and min-max strategies are called \textbf{optimal strategies}.
\end{defbox}

Suppose $s_1^*$ is a max-min strategy for P1, and $s_2^*$ is a min-max strategy for P2.
We always have:
\[
  \underline{v} \leq u(s_1^*, s_2^*) \leq \overline{v}
\]
If the game has a value ($\underline{v} = \overline{v} = v$), then $u(s_1^*, s_2^*) = v$.
Furthermore:
\[
  u(s_1, s_2^*) \leq v \quad \text{and} \quad u(s_1^*, s_2) \geq v
\]

% ============================================================
\section{Nash Equilibrium Connection}

\begin{defbox}{Nash Equilibrium (NE)}
  A strategy profile $(s_1^*, \dots, s_n^*)$ is called a \textbf{Nash Equilibrium} if no player has an incentive to unilaterally deviate. For all $i$:
  \[
    u_i(s_i^*, s_{-i}^*) \geq u_i(s_i, s_{-i}^*) \quad \forall s_i \in S_i
  \]
\end{defbox}

\begin{thmbox}{Minimax Theorem (Partial Statement)}
  For a two-player zero-sum game, if the value exists ($\underline{v} = \overline{v}$), then $(s_1^*, s_2^*)$ is a Nash Equilibrium, where $s_1^*$ is a max-min strategy of Player 1, and $s_2^*$ is a min-max strategy for Player 2.
\end{thmbox}

\begin{takebox}
  In zero-sum games with a defined value, playing defensively (optimizing your worst-case scenario) coincides perfectly with finding the Nash Equilibrium. Your best defensive move is also your best offensive response to an optimal opponent!
\end{takebox}

\end{document}